%قبل از شروع مسئله ترجیحاً چیزی ننویسید. اما اگر خواستید بنویسید، در قالب لیستی (فقرات و شمارش) بنویسید.
%برای اضافه کردن نمره، بعدی از شروع‌مسئله، [x] را اضافه کنید و به جای x  نمره مورد نظر را قرار دهید. برای زیرمسئله‌ها هم به همین صورت است.
%\شروع‌مسئله[10]

\شروع‌مسئله[16]
$n$ 
نقطه روی صفحه مختصات به شنگدباو داده‌شده‌است. او می‌خواهد تعداد مربع‌هایی را بیابد که رئوس آن روی این نقاط باشند.
برای این‌کار الگوریتم زیر را پیشنهاد کرده‌است:

هر سری بین تمامی خطوط افقی و عمودی که روی آن حداقل یک نقطه باشد، آن خطی را در نظر می‌گیریم که تعداد نقاط رویش کمینه‌باشد.
سپس به ازای هر جفت از نقاط که روی آن خط هستند چک می‌کنیم آیا مربعی هست که این دو نقطه دو سر یکی از اضلاعش باشند.
توجه کنید حداکثر دو مربع می‌توانیم داشته‌باشیم که دو نقطه‌اش نقاطی باشند که ما در نظر گرفته‌ایم.
پس از اینکه تمامی مربع‌هایی را یافتیم که نقاط رئوسش روی این خط هستند. تمامی نقاط روی این خط را پاک می‌کنیم و به صورت بازگشتی دو‌باره الگوریتم را تکرار می‌کنیم.

فرض کنید چک کردن اینکه ۴ رأس آیا باهم تشکیل مربع می‌دهند یا نه ۱ واحد هزینه‌ دارد. از این نکته استفاده کنید که خطی افقی یا عمودی وجود دارد که روی آن حداکثر 
$\sqrt{n}$
تا نقطه‌باشد؛ سپس با استفاده از این نکته ثابت کنید 
$O(n\sqrt{n})$
واحد هزینه کافیست بپردازیم. سپس از این موضوع نتیجه بگیرید که تعداد مربع‌هاییکه رئوس آن روی
این نقاط باشند از همین مرتبه‌است.

%اگر مسئله چند قسمت داشت به این صورت عمل کنید:
%\شروع{زیرمسائل}
%\زیرمسئله[0] % مسئلۀ ۱آ
%\زیرمسئله % مسئلۀ ۱ب
%\زیرمسئله % مسئلۀ ۱پ
%\زیرمسئله % مسئلۀ ۱ت
%\پایان{زیرمسائل}

\solution{
	فرض کنید تعداد
	$a$
	تا نقطه پیدا کرده باشیم که با هم روی یک خط (عمودی یا افقی) قرار گرفته باشند. باتوجه به اینکه مربع ها افقی یا عمودی هستند مربع‌هاییکه ممکن است از روی این نقاط ساخته شوند
	$a \times (a-1)$
	تا هستند. یعنی به ازای هر جفت نقطه می‌توانیم مربع‌هاییکه این جفت یک ضلعش را تشکیل دهند چک کنیم. کافیست ببینیم در نقاط متناظر در دو طرف خط آیا مربعی وجود دارد یا خیر.
	حالا که همه مربع‌هاییکه با این نقاط می‌تواند به وجود بیاید را پیدا کردیم می‌توانیم آن نقاط را حذف کنیم و بقیه مسئله را مشابهاً حل کنیم.
	بنابراین اگر خطی افقی یا عمودی پیدا کنیم که روی آن 
    مجموعه 
    $|A|$
    از نقاط باشد. و $S$
    مجموعه همه نقاط‌مان باشد و 
    $T(n)$
    برابر تعداد مربع‌های تشکیل شده از این نقاط باشد داریم:
    
    $$T(S) = T(S-A) + a\times(a-1)$$
    
    حال ثابت می‌کنیم مجموعه $A$
    از نقاط وجود دارد که
    $|A| \le \sqrt{|S|}$
    
    برای اثبات این موضوع فرض کنید همه خطوط افقی لااقل
    $\sqrt{|S|}$
    تا نقطه دارند. (توجه کنید منظورمان بین خطوطی هست که لااقل یک نقطه رویشان ظاهر شده‌باشد)
    اگر اینطور نباشد یکی از آن خطوط را با نقاط رویش در نظر بگیریم به ما مجموعه 
    $A$
    مورد نظر را خواهد داد.
    اما اگر همه از این مقدار بیشتر باشند واضح است که تعداد خطوط افقی‌ای که لااقل 
    $\sqrt{|S|}$
    تا نقطه دارند حداکثر
    $\sqrt{|S|}$
    تا اند! (چرا؟)
    
    بنابراین هر خط عمودی‌ای که در نظر بگیریم که حداقل یک نقطه رویش باشد. در نقاط تقاطعش با خطوط افقی ممکن است نقطه دیگری داشته باشند.
    بنابراین چون می‌دانیم تعداد خطوط افقی در این حالت حداکثر
    $\sqrt{|S|}$
    تاست می‌توان دید که 
    تعداد نقاط روی این خط عمودی نیز حداکثر این تعداد است.
    
    پس همیشه می‌توان 
    $A$
    ای پیدا کرد که 
    $|A| \le \sqrt{|S|}$
    حالا رابطه بازگشتی 
    $$T(S) = T(S-A) + (|A|) \times (|A| - 1)$$
    را اگر باز کنیم به صورت مجموع مربعات یکسری عدد در می‌آید که مجموع خودشان برابر
    $|S| = n$
    است. می‌دانیم در بدترین حالت این اعداد 
    به صورت یکسری اعداد
    $\sqrt{n}$
    هستند (برای اثبات آن می‌توانیم دو تا مجموعه که اندازه شان کمتر از این مقدار است را در نظر بگیریم با انتقال یک عضو از مجموعه کوچکتر به بزرگتر مجموع ثابت می‌ماند ولی مجموع مربعات زیاد می‌شود)
    پس در بدترین حالت
    جمع 
    $\sqrt{n}$
    تا عدد برابر
    $\sqrt{n}$
    است که مربع‌هایشان برابر
    $n$
    می‌شود یعنی در کل مجموع مربعات حداکثر 
    $n\sqrt{n}$
    است.
    
    }




%%%%%%%%%%%%%%%حداکثر بخط وجود دارد.ه همین تعداد نقطه روی این %%%%%%%%%%%%%%%%%%%%%%%%%%%%%%%%%%%%%%
% برخی دستورات مفید برای تمرین‌ها. %
%%%%%%%%%%%%%%%%%%%%%%%%%%%%%%%%%%%%%%%%%%%%%%%%%%%%%

% فرمول‌های نمونه:
%$f(x) = \Theta(x)$
%$T(x, y) \leq \log(x) + 2^y + \binom{2n}{n}$
% {\tt code\_function}


% لیست موارد بدون شماره‌گذاری:
%\شروع{فقرات}
%    \فقره مورد اول 
%        \شروع{فقرات}
%            \فقره مورد اول داخلی 
%                \شروع{فقرات}
%                    \فقره اولین مورد خیلی داخلی 
%                    \فقره دومین مورد خیلی داخلی 
%                \پایان{فقرات}
%            \فقره مورد دوم داخلی 
%        \پایان{فقرات}
%    \فقره مورد دوم 
%\پایان{فقرات}

% لیست موارد شماره‌گذاری شده:
%\شروع{شمارش}
%    \فقره یک مورد 
%    \فقره یک مورد دیگر 
%    \فقره باز هم یک مورد دیگر
%\پایان{شمارش}

% برای نوشتن یک فرمول چند خطی که یک شماره بخوره. برای حالت‌های دیگر به سایت زیر مراجعه کنید:
% https://www.overleaf.com/learn/latex/Aligning_equations_with_amsmath
%\begin{align} 
%    \begin{split}
%        (x+y)^3 &= (x+y)^2(x+y) \\
%                &= (x^2+2xy+y^2)(x+y) \\
%                &= (x^3+2x^2y+xy^2) + (x^2y+2xy^2+y^3) \\
%                &= x^3+3x^2y+3xy^2+y^3
%    \end{split}                                 
%\end{align}

% نوشتن ماتریس و آرایه دوبعدی
%\begin{align}
%    A = 
%    \begin{bmatrix}
%        A_{11} & A_{21} \\
%        A_{21} & A_{22}
%    \end{bmatrix}
%\end{align}

% اضافه کردن تصویر با شرح. برای اضافه کردن تصویر بدون شرح، فقط دستور درج‌تصویر را استفاده کنید.
%\شروع{شکل}[h!]
%\تنظیم‌ازوسط
%  \درج‌تصویر[پهنا=0.5\پهنای‌متن]{img.jpg}
%\شرح{شرح تصویر}
%\برچسب{تصویر:۱}
%\پایان{شکل}
%\پایان{زیرمسائل}

%کد زیر را ببینید.
%%\lstinputlisting{source_filename.py}
%\begin{latin}
%\begin{lstlisting}
%# Import the modules
%import sys
%import random
%
%ans = True
%
%while ans:
%question = raw_input("Ask the magic 8 ball a question: (press enter to quit) ")
%
%answers = random.randint(1,8)
%
%if question == "":
%sys.exit()
%
%elif answers == 1:
%print "It is certain"
%
%elif answers == 2:
%print "Outlook good"
%
%elif answers == 3:
%print "You may rely on it"
%
%elif answers == 4:
%print "Ask again later"
%
%elif answers == 5:
%print "Concentrate and ask again"
%
%elif answers == 6:
%print "Reply hazy, try again"
%
%elif answers == 7:
%print "My reply is no"
%
%elif answers == 8:
%print "My sources say no"
%\end{lstlisting}
%\end{latin}
